\section{An inverse estimate for a witnessing edge}
\label{sec:inverse}

The exceptional compact nodes will be controlled by varying the isolated
prime in a regular box.  The first step identifies the only obstruction
to this variation: a positive proportion of the relevant primes must
share a rational intercept. The use of modular concentration to obtain
algebraic structure is related to the inverse-sieve methods of
\cite{HelfgottVenkatesh2009,Walsh2012}. More specifically, the
small-height polynomial and prime-product strategy is related to Walsh's
algebraicity method \cite[Section~3]{Walsh2014}. The partial-incidence
modular-line form needed in a prime bin, including its uniformity across
the independent boxes, is proved here rather than invoked from that work.

\begin{definition}
A prime \(p\) \emph{aligns} with a reduced rational \(A/D\), where
\(A\in\Z\), \(D\in\N\), and \(\gcd(A,D)=1\), if
\[
  D a_p\equiv A\pmod p.
\]
This condition implies \(p\nmid D\).
For a squarefree product \(q\) of primes carrying prescribed classes,
write
\[
 q_E(A/D)=
 \prod_{\substack{\ell\mid q\\
                   \ell\text{ does not align with }A/D}}\ell .
\]
\end{definition}

Throughout this section put
\begin{equation}
 \mathcal Z=\exp\!\left(\frac{L}{\log L}\right).
 \label{eq:inverse-z}
\end{equation}
As usual, all parameters other than \(z\) are fixed before taking a limit
in \(z\).  In particular, a statement valid for every sufficiently small
fixed \(\xi>0\) allows its eventual lower threshold on \(z\) to depend on
\(\xi\).

Fix one candidate edge in a full regular box from
\cref{prop:regular-boxes}.  Write its parent as \(p q\), with \(p\) the
isolated prime, and denote its edge prime by \(u\).  Let
\(\mathcal P,\mathcal U\) be their respective prime bins, let \(R,U\) be
the lower endpoints of these bins, and put
\[
 n_p=\#\mathcal P,\qquad n_u=\#\mathcal U.
\]
Both bins are right-closed and have upper-to-lower endpoint ratio at most
\(1+\xi\).  Their logarithmic widths are comparable to \(\xi\); hence,
for each fixed \(\xi>0\),
\begin{equation}
 n_p\asymp \frac{\xi R}{\log R},
 \qquad
 n_u\asymp \frac{\xi U}{\log U}.
 \label{eq:inverse-bin-counts}
\end{equation}
The constants here can be chosen uniformly for \(0<\xi\le1\), once \(z\)
is sufficiently large depending on \(\xi\).

Let \(\mathcal Q\) be the set of possible products \(q\).  Such a product
uniquely determines its sorted prime choices.  Denote by \(S_q\) the
product of the upper endpoints of its specified bins, with the
prescribed multiplicities; \(S_q\) is therefore fixed for the box.
The hypotheses supplied by \cref{prop:regular-boxes} include
\begin{align}
 &z^{\alpha_{\min}}\le R\le z^{\alpha_{\max}},
       \qquad \alpha_{\max}=\frac1{100},
       \qquad w\le U\le w^{2K_*+4},                                      \label{eq:inverse-box-ranges}\\
 &\log_w\frac{Y}{p q}\le2K_*+4,
       \qquad
       \log_w\frac{Y}{p q u}\ge\frac{3a_\star}{4},                       \label{eq:inverse-length-ranges}\\
 &1\le \frac{S_q}{q}\le
       (1+\xi)^{C_{\rm len}\log B}.                                     \label{eq:inverse-q-range}
\end{align}
The bins of \(p\) and \(u\) each have multiplicity one in the full box.
Thus the choices of \(p,q,u\), and all descendants after the edge are
independent in the full box, and \(p,u\) divide none of the other
products involved.  In particular,
\[
 \gcd(p,q)=\gcd(pq,u)=1 .
\]
Writing \(\alpha=\log R/L\), the parent-length bounds imply
\begin{equation}
 q=z^{\,2-\alpha+o(1)},\qquad
 S_q=z^{\,2-\alpha+o(1)},
 \qquad
 \alpha_{\min}\le\alpha\le\alpha_{\max}.
 \label{eq:inverse-q-size}
\end{equation}
All the \(o(1)\)'s in this section are uniform over boxes satisfying
these fixed-parameter bounds, for each permitted fixed \(\xi\).

The normalized \emph{harmonic measure} on the box assigns mass
proportional to the reciprocal of the product of its chosen primes.
On \(p,q,u\) this is the product measure with respective weights
\(1/p,1/q,1/u\).  Descendant choices may be integrated out whenever an
event depends only on \(p,q,u\).

\begin{proposition}[Inverse estimate for an edge]
\label{prop:inverse}
Fix the parameters in \eqref{eq:inverse-box-ranges}--%
\eqref{eq:inverse-q-range}, the witnessing threshold \(\delta>0\), and
\(0<\sigma<1\).
There are constants \(c_p>0\), a list-size bound \(F<\infty\), and
\(\xi_{\rm inv}>0\), independent of the eventual choice of fixed
\(0<\xi\le\xi_{\rm inv}\), with the following property for all
sufficiently large \(z\).

The full candidate box has a list \(\mathcal A\) of at most \(F\)
reduced rationals \(A/D\) satisfying
\begin{equation}
 D\le R\mathcal Z^{10},\qquad
 |A|\le S_qR\mathcal Z^{10},\qquad
 \#\{p\in\mathcal P:p\text{ aligns with }A/D\}\ge c_p n_p .
 \label{eq:inverse-list}
\end{equation}
Outside a set of normalized harmonic box measure at most \(\sigma\),
every occurrence of
\begin{equation}
 \left|N_{pqu}-\frac{N_{pq}}u\right|
       >\delta\,\frac{YV_0}{pqu}
 \label{eq:inverse-bad-edge}
\end{equation}
has a rational \(A/D\in\mathcal A\) for which
\begin{equation}
 p\text{ aligns with }A/D,
 \qquad D q_E(A/D)\le R\mathcal Z^{10}.
 \label{eq:inverse-structure}
\end{equation}
The exceptional measure is relative to the whole box, rather than to
the subset on which \eqref{eq:inverse-bad-edge} holds.
\end{proposition}

The proof fixes the whole-box rational list before arguing by
contradiction. If too many bad edges fail \eqref{eq:inverse-structure},
sampling their frozen residue tests produces many parent pairs but few
permissible slopes. An integer polynomial with small coefficient height
then vanishes on the corresponding integer points. The lemma below bounds
the contribution of its nonlinear components. A remaining rational line
has an intercept in the fixed list, forcing most selected pairs on that
line to satisfy \eqref{eq:inverse-structure} after all.

The lattice-point estimate below combines algebraic decomposition with
Jarn\'{i}k's convex-arc argument \cite[Sections~1--2]{Jarnik1926}.
We give the proof with explicit degree dependence because the degree
grows with the prime bin. The coefficient-independent intersection
bound used in its proof is established in
Lemma~\ref{lem:elementary-bezout}.

\begin{lemma}[Degree-uniform nonlinear lattice-point bound]
\label{lem:nonlinear-points}
Let \(P\in\R[X_1,X_2]\) be nonzero and of degree at most \(D_0\ge1\),
and let \(S\ge1\).
The number of lattice points in \([0,S]^2\) lying on at least one
nonlinear irreducible component of \(P=0\), with irreducibility taken
over \(\mathbb C\), is
\[
 O\!\left(D_0^4(1+S^{2/3})\right).
\]
The implied constant is absolute, independently of the coefficients
of \(P\).  No bound is asserted for points lying only on line
components.
\end{lemma}

\begin{proof}
Consider first one nonlinear irreducible factor \(Q\) of degree \(d\).
If \(Q\) is not proportional to a real polynomial, \(Q\) and its
coefficientwise conjugate are distinct irreducible polynomials.  Every
real zero of \(Q\) is a common zero, so B\'ezout's theorem gives at most
\(d^2\) such points.

Suppose instead that \(Q\) has been scaled to have real coefficients.
Neither \(Q_{X_1}\) nor \(Q_{X_2}\) vanishes identically: otherwise
irreducibility over \(\mathbb C\) would force \(Q\) to be linear.
Their degrees are smaller than \(d\), so their intersections with
\(Q=0\) give \(O(d^2)\) points in total.  These include all singularities
and all horizontal or vertical tangencies.
There are also \(O(d)\) intersections with the boundary of the square.

On a smooth graph \(X_2=f(X_1)\) with \(Q_{X_2}\ne0\), the second
derivative vanishes precisely when
\begin{equation}
 \mathcal I_Q
 =Q_{X_1X_1}Q_{X_2}^{\,2}
  -2Q_{X_1X_2}Q_{X_1}Q_{X_2}
  +Q_{X_2X_2}Q_{X_1}^{\,2}
 =0.
 \label{eq:inverse-inflection}
\end{equation}
Indeed \(f''=-\mathcal I_Q/Q_{X_2}^3\) on the graph.
If \(Q\nmid\mathcal I_Q\), B\'ezout's theorem adds \(O(d^2)\)
inflection points, because \(\deg\mathcal I_Q\le3d-4\).
If \(Q\mid\mathcal I_Q\), any nonspecial interior smooth graph arc has
\(f''=0\) throughout and therefore lies on a line.  Vanishing on that
arc would make the line a factor of \(Q\), a contradiction.
Thus in this latter case all real points in the square are already
among the critical or boundary points counted above.

In the remaining case, mark the abscissae of all critical, boundary,
and inflection points, together with \(0,S\).
There are \(O(d^2)\) marked abscissae.
Between consecutive marked abscissae the interior real zero set is a
union of at most \(d\) smooth graphs extending across that open
interval.  To justify extension, the implicit function theorem applies
because \(Q_{X_2}\ne0\); continuation cannot end inside the interval
without a critical point or an intersection with the boundary.
The bound \(d\) follows from the degree of a vertical fiber.
Each graph is monotone, and its derivative is strictly monotone,
because neither \(Q_{X_1}\) nor \(\mathcal I_Q\) vanishes on it.

Consider the lattice points of one such graph, ordered by increasing
abscissa.  Consecutive difference vectors have positive integer first
coordinate and distinct slopes, by strict convexity or strict
concavity.  They are therefore distinct integer vectors.
Their total \(\ell^1\)-length is at most \(2S\), since both graph
coordinates are monotone and range over intervals of length at most
\(S\).  There are \(O(T^2)\) integer vectors of \(\ell^1\)-length at most
\(T\).  Consequently \(m\) distinct such vectors have total length
\(\gg m^{3/2}\): take \(T\) to be a sufficiently small fixed multiple
of \(m^{1/2}\), so that at least half of the vectors have length greater
than \(T\).  It follows that this graph contains
\(O(1+S^{2/3})\) lattice points.

The marked vertical fibers contain at most \(d\) points each, since
an identically vanishing fiber would give a vertical line factor of
\(Q\).  Summing over the \(O(d^3)\) graphs and the marked fibers proves
the bound \(O(d^3(1+S^{2/3}))\) for this component.
Finally, sum over the distinct nonlinear factors of \(P\).  Their
degrees sum to at most \(D_0\), so the asserted, slightly looser,
bound with \(D_0^4\) follows.
\end{proof}

\begin{proof}[Proof of \cref{prop:inverse}]
\textbf{The fixed whole-box list.}
To make the order of choices explicit, choose an integer \(h_1\) such
that
\[
 \frac{a_\star h_1}{2}>12,
\]
and set
\begin{equation}
 \beta=\frac{\sigma}{16},\qquad
 \gamma_{\rm inc}=\frac{\beta}{2}
       \left(\frac{\beta}{4}\right)^{h_1},\qquad
 c_5=\frac{\gamma_{\rm inc}}{8},\qquad
 c_p=\frac{\gamma_{\rm inc}}{32}.
 \label{eq:inverse-fixed-constants}
\end{equation}
These constants are fixed before \(\xi\) is chosen and before any
failure set, cell, or sample is selected.
For each candidate box let \(\mathcal A\) be the set of \emph{all}
reduced rationals satisfying \eqref{eq:inverse-list}.
This is a finite set by the height restrictions.

If two distinct list rationals \(A_1/D_1,A_2/D_2\) align at a prime,
that prime divides the nonzero determinant \(A_1D_2-A_2D_1\).
Its absolute value is at most
\[
 2 S_q R^2\mathcal Z^{20}=z^{\,2+\alpha+o(1)}.
\]
Every prime of \(\mathcal P\) is at least \(z^{\alpha_{\min}}\).
Thus two distinct rationals have at most \(C_{\alpha_{\min}}\)
common aligning primes, with a fixed constant independent of \(\xi\).
If \(M=\#\mathcal A\), and \(m_p\) counts the list rationals aligning
at \(p\), Cauchy's inequality gives
\[
 M^2c_p^2 n_p
 \le\sum_{p\in\mathcal P}m_p^2
 \le M n_p+C_{\alpha_{\min}}M(M-1).
\]
Once \(n_p\ge2C_{\alpha_{\min}}/c_p^2\), this implies
\[
 M\le \frac{2}{c_p^2}.
\]
We may therefore take \(F=\lceil2/c_p^2\rceil\).
This bound is independent of \(\xi\); the threshold on \(z\) may depend
on \(\xi\) through \(n_p\).

Call a pair \((p,q)\) \emph{structured} if
\eqref{eq:inverse-structure} holds for some member of this fixed
whole-box list.  A triple \((p,q,u)\) is called \emph{bad} if its parent
pair is nonstructured and \eqref{eq:inverse-bad-edge} holds.
Assume, towards a contradiction, that bad triples occupy more than a
\(\sigma\) fraction of the full harmonic box measure. Descendants can
now be integrated out.

\medskip
\textbf{From bad triples to sampled incidences.}
We will select a fixed number of edge primes so that many nonstructured
parent pairs satisfy all their frozen residue tests, while the total
number of permitted slopes over the isolated primes is small.
First we restrict \(q\) to a
sufficiently large cell. For \(q\in\mathcal Q\), let
\(b_q\in[1,q]\cap\Z\) be its
Chinese remainder theorem (CRT) representative for the prescribed hits.  Put
\[
 M_0=\prod_{t\le w}t ,
\]
where the product is over primes.  By the prime number theorem,
\[
 \log M_0=O(w)
 =O\!\left(\frac{L}{(\log L)^2}\right)=o(\log\mathcal Z).
\]
In particular \(M_0^2=\mathcal Z^{o(1)}\).
Partition \(\mathcal Q\) into cells by a value interval of ratio at
most \(1+\xi\), and by the pair
\((q,b_q)\bmod M_0\).  By \eqref{eq:inverse-q-range} the number of value
intervals needed is \(O(C_{\rm len}\log B+1)\), so the total number of
cells is \(\mathcal Z^{o(1)}\).

We record the cardinality bound
\begin{equation}
 \#\mathcal Q\ge \frac{S_q}{\mathcal Z^{o(1)}}.
 \label{eq:inverse-many-q}
\end{equation}
Indeed, a bin with upper endpoint \(P_j\) and multiplicity \(m_j\)
contains \(\gg_\xi P_j/\log P_j\) primes.  Even in the smallest bin
this number is much larger than \(C_{\rm len}\log B\), for fixed
\(\xi>0\) and large \(z\).
If \(n_j\) is its prime count and
\(k=\sum_jm_j\le C_{\rm len}\log B=O(\log L)\), then
\[
 \prod_j\binom{n_j}{m_j}
 \ge \prod_j\frac{(n_j/2)^{m_j}}{m_j!}
 \ge S_q\exp\!\bigl(-O_\xi((\log L)^2)\bigr).
\]
The first inequality holds once \(n_j\ge2m_j\).
In the second, the losses from \(\log P_j\), fixed bin-width constants,
and the factorials have total logarithm \(O_\xi((\log L)^2)\).
Unique factorization identifies these choices with distinct \(q\)'s.
Since \((\log L)^2=o(\log\mathcal Z)\), this proves
\eqref{eq:inverse-many-q}.

Discard cells containing fewer than \(S_q/\mathcal Z\) points.
They contain at most \(S_q\mathcal Z^{-1+o(1)}\) points in total.
The weights \(1/q\) vary globally by at most
\((1+\xi)^k=\mathcal Z^{o(1)}\), so their total normalized harmonic
mass is \(\mathcal Z^{-1+o(1)}=o(1)\).
For sufficiently large \(z\), a retained initial cell \(\mathcal C\)
therefore has conditional harmonic density of bad triples at least
\(\sigma/2\), and
\begin{equation}
 \#\mathcal C\ge S_q/\mathcal Z.
 \label{eq:inverse-cell-size}
\end{equation}
Within \(\mathcal P,\mathcal C,\mathcal U\) each reciprocal weight
varies by a factor at most \(1+\xi\).
As long as \(\xi\le1\), converting their product measure to uniform
counting costs at most \((1+\xi)^3\le8\).
The count density of bad triples in
\(\mathcal P\times\mathcal C\times\mathcal U\) is consequently at least
\(\beta=\sigma/16\), with \(\beta\) fixed in
\eqref{eq:inverse-fixed-constants}.

\medskip
\emph{Freezing the parent interval.}
We now place the small-prime count on an interval that is fixed as
\(q\) varies within \(\mathcal C\).
For fixed \(p\in\mathcal P\) and \(q\in\mathcal C\), put
\begin{equation}
 s_p(q)=(a_p-b_q)\overline q\pmod p,
 \qquad 0\le s_p(q)<p.
 \label{eq:inverse-s}
\end{equation}
The parent hit progression is
\[
 n=b_q+q(s_p(q)+pj),\qquad j\ge0.
\]
The initial hit belongs to \([1,pq]\), so the condition \(j\ge0\)
indexes exactly the positive hits.
The allowed \(j\)'s form an initial interval of integers whose length
is \(Y/(pq)+O(1)\).

If the value range of \(\mathcal C\) is contained in
\([Q_-,Q_+]\), with \(Q_+/Q_-\le1+\xi\), choose a common initial
integer interval \(I_p\) of length
\[
 J=J_p=\left\lfloor\frac{Y}{pQ_+}\right\rfloor.
\]
Its symmetric difference with the actual interval has length
\(O(\xi J+1)\).  The classes to avoid at primes \(t\le w\) on \(I_p\)
depend only on \(p\), the fixed cell residues, and the value
\(v=s_p(q)\bmod M_0\).  This uses the invertibility of \(pq\) at every
such prime.  For each possible \(v\), let \(m=m_{p,v}\) count the
survivors on \(I_p\) for that pattern, and let \(m_{e,u}=m_{p,v;e,u}\)
count those with \(j\equiv e\pmod u\).

The upper bounds in \cref{lem:small-sieve} give, uniformly in the
patterns,
\begin{equation}
 m\ll JV_0.
 \label{eq:inverse-m-upper}
\end{equation}
They also show that replacing intervals changes the survivor count
by \(O(\xi JV_0)\), and the count in any specified \(j\)-class modulo
\(u\) by \(O(\xi JV_0/u)\).
For the latter assertion, enclose the difference subsequence in an
interval of length comparable to \(\xi J/u\).
By \eqref{eq:inverse-length-ranges}, for fixed \(\xi>0\) this length
is eventually much larger than \(w^{a_\star/2}\); the corresponding
progression step \(pqu\) is a unit at every prime at most \(w\).
The same reasoning applies to the unconditioned difference interval.
The upper-sieve constants are independent of fixed \(\xi\), although
the threshold at which these lengths suffice may depend on \(\xi\).

The edge hit selects the class
\begin{equation}
 e_u=e_u(p,q)=(a_u-b_q-q s_p(q))\overline{pq}\pmod u .
 \label{eq:inverse-edge-class}
\end{equation}
Choose \(\xi_{\rm inv}\le1\) sufficiently small, depending on
\(\delta\) and the fixed upper-sieve constants.
The preceding replacement estimates show that
\eqref{eq:inverse-bad-edge} implies
\begin{equation}
 |m_{e_u,u}-m/u|>\delta'\frac{JV_0}{u},
 \label{eq:inverse-frozen-deviation}
\end{equation}
for a fixed \(\delta'>0\).
For example, after decreasing \(\xi_{\rm inv}\) the choice
\(\delta'=\delta/4\) suffices.
Define \(E_{p,u}(v)\) to be the set of residues \(e\bmod u\) satisfying
\eqref{eq:inverse-frozen-deviation}.
These sets depend on \(p\), the initial cell, and \(v\), and no
independence in those variables is asserted.

\medskip
\emph{The large-sieve saving.}
We bound the average proportion of residues passing an exceptional test.
For fixed \(p,v\), let \(\mathcal S\subset I_p\) be the surviving
sequence, so \(\#\mathcal S=m\), and put
\[
 S(\theta)=\sum_{j\in\mathcal S}\exp(2\pi i j\theta).
 \]
Additive orthogonality gives the exact identity
\begin{equation}
 u\sum_{e\bmod u}(m_{e,u}-m/u)^2
   =\sum_{a=1}^{u-1}|S(a/u)|^2.
 \label{eq:inverse-orthogonality}
\end{equation}
For prime \(u\), all fractions \(a/u\) on the right are reduced.
Fractions belonging to distinct such primes are distinct, and the
whole family is separated modulo one by at least
\(((1+\xi)U)^{-2}\ge(4U^2)^{-1}\).
Lemma~\ref{lem:elementary-large-sieve} therefore gives
\[
 \sum_{u\in\mathcal U}
 u\sum_{e\bmod u}(m_{e,u}-m/u)^2
 \ll (J+U^2)m;
\]
see \cite{MontgomeryVaughan1973,VaughanLargeSieve}.
No condition other than interval support is imposed on the sequence
\(\mathcal S\).
Each exceptional residue contributes more than
\((\delta'JV_0)^2/u\) to the corresponding left side.
Using \eqref{eq:inverse-m-upper}, we obtain
\begin{align}
 \frac1{n_u}\sum_{u\in\mathcal U}\frac{|E_{p,u}(v)|}{u}
 &\ll
 \frac{(J+U^2)m}{n_u(JV_0)^2} \notag\\
 &\ll_\xi(\log w)^2\left(\frac1U+\frac UJ\right)
 \le w^{-a_\star/2}.
 \label{eq:inverse-large-sieve}
\end{align}
For the last inequality, \(\log U=O_{K_*}(\log w)\),
\(V_0\asymp1/\log w\), \(U\ge w\), and
\(J/U\gg w^{3a_\star/4}\).
The fixed factor depending on \(\xi\), as well as the logarithmic
factors, is absorbed between the exponents \(3a_\star/4\) and
\(a_\star/2\) by taking \(z\) large.
This proves \eqref{eq:inverse-large-sieve} uniformly in \(p,v\) and
the chosen cell, in exactly the fixed-\(\xi\) sense required here.

\medskip
\emph{Selecting the sampled relation.}
For \((p,q)\in\mathcal P\times\mathcal C\), let \(f(p,q)\) be the
fraction of \(u\in\mathcal U\) for which \((p,q,u)\) is bad.
Its uniform count average is at least \(\beta\).
At least a \(\beta/2\) fraction of pairs therefore have
\(f(p,q)\ge\beta/2\): the complementary bound follows at once from
\(0\le f\le1\).

Sample an ordered \(h_1\)-tuple of distinct primes
\(\mathbf u=(u_1,\ldots,u_{h_1})\) uniformly from \(\mathcal U\).
For this sample define the pair relation
\[
\begin{split}
 \mathcal I_{\mathbf u}=\bigl\{(p,q)\in\mathcal P\times\mathcal C:
 &\ (p,q)\text{ is nonstructured, and}\\
 &\ e_{u_i}(p,q)\in E_{p,u_i}(s_p(q)\bmod M_0)
       \text{ for every }1\le i\le h_1\bigr\},
\end{split}
\]
and put \(I(\mathbf u)=\#\mathcal I_{\mathbf u}/(n_p\#\mathcal C)\).
Members of \(\mathcal I_{\mathbf u}\) will be called \emph{sampled
incidences}. By \eqref{eq:inverse-frozen-deviation}, a pair for which
every sampled triple is bad belongs to \(\mathcal I_{\mathbf u}\).
This implication suffices for the lower bound on its expected density.
Once \(n_u\ge4h_1/\beta\), each pair with \(f(p,q)\ge\beta/2\) has probability
at least \((\beta/4)^{h_1}\) that all sampled primes witness
\eqref{eq:inverse-bad-edge}.
Consequently
\begin{equation}
 \E I(\mathbf u)\ge
 \frac{\beta}{2}\left(\frac{\beta}{4}\right)^{h_1}
       =\gamma_{\rm inc}.
 \label{eq:inverse-incidence-expectation}
\end{equation}

Write \(W=\prod_{i=1}^{h_1}u_i\) and refine \(\mathcal C\) by fixing
\((q,b_q)\bmod W\).  For any such refinement \(\mathcal C'\), and
fixed \(p\), let \(\mathcal T_p(\mathcal C')\subset[0,p)\cap\Z\)
consist of those integers \(s\) satisfying all the sampled tests,
using \(v=s\bmod M_0\).
For a fixed \(v\), \eqref{eq:inverse-edge-class} is an invertible
affine function of \(s\bmod u_i\), with coefficient
\(-\overline p\bmod u_i\).
Its other coefficients are fixed by the refinement.
Thus it restricts \(s\bmod u_i\) to exactly
\(|E_{p,u_i}(v)|\) residues.
CRT, applied also to \(s\equiv v\pmod {M_0}\), gives
\begin{align}
 |\mathcal T_p(\mathcal C')|
 &\le K_p(\mathbf u),\notag\\
 K_p(\mathbf u)
 &:=
 \frac{p}{M_0}\sum_{v\bmod M_0}
       \prod_{i=1}^{h_1}\frac{|E_{p,u_i}(v)|}{u_i}
       +M_0 W .
 \label{eq:inverse-crt-slopes}
\end{align}
Indeed, for each \(v\) the number of allowed residues modulo \(M_0W\)
is \(\prod_i|E_{p,u_i}(v)|\); an interval of length \(p\) contains at
most \(p/(M_0W)+1\) integers in each residue.
Summing the rounding terms over \(v\) gives at most \(M_0W\).
All CRT factors are coprime, and every sampled \(u_i\) is coprime to
\(pq\) by the singleton-bin hypotheses.
Crucially, the upper bound \(K_p(\mathbf u)\) is independent of the
refinement centers.

For any numbers \(0\le t_u\le1\), their product averaged over distinct
ordered \(h_1\)-tuples is at most
\[
 \frac{n_u^{h_1}}{(n_u)_{h_1}}
 \left(\frac1{n_u}\sum_{u\in\mathcal U}t_u\right)^{h_1},
 \qquad
 (n_u)_{h_1}=n_u(n_u-1)\cdots(n_u-h_1+1).
\]
The prefactor tends to one since \(h_1\) is fixed.
Applying \eqref{eq:inverse-large-sieve} for each \(p,v\) in
\eqref{eq:inverse-crt-slopes} yields
\begin{equation}
 \E\frac{K_p(\mathbf u)}p
 \le 2w^{-a_\star h_1/2}+z^{-c'}
 \label{eq:inverse-slopes-expectation}
\end{equation}
for some fixed \(c'>0\).
Here \(M_0W=z^{o(1)}\), since \(h_1\) and the edge-bin exponent range are
fixed, whereas \(p\ge z^{\alpha_{\min}}\).
Since \(p/R\le1+\xi\le2\), the same estimates imply
\[
 \E\left(\frac1{n_pR}\sum_{p\in\mathcal P}K_p(\mathbf u)\right)
 \ll w^{-a_\star h_1/2}+z^{-c'}
 =o(L^{-6}).
\]
The last step uses \(w=L/(\log L)^2\) and \(a_\star h_1/2>12\).
By Markov's inequality, the probability that
\(\sum_pK_p>n_pR/L^6\) is \(o(1)\).
As \(0\le I(\mathbf u)\le1\), removing those samples changes the
expectation in \eqref{eq:inverse-incidence-expectation} by \(o(1)\).
Hence there exists a sample for which
\begin{equation}
 I(\mathbf u)\ge\gamma_{\rm inc}/2,
 \qquad
 \sum_{p\in\mathcal P}K_p(\mathbf u)\le\frac{n_pR}{L^6}.
 \label{eq:inverse-good-sample}
\end{equation}
Fix this sample and abbreviate its bounds by \(K_p\).

There are at most \(W^2=w^{O(1)}=\mathcal Z^{o(1)}\) refined cells.
Discard those with fewer than \(S_q/\mathcal Z^3\) points.
Relative to \(\mathcal C\), whose size satisfies
\eqref{eq:inverse-cell-size}, their total count fraction is at most
\(W^2/\mathcal Z^2=o(1)\).
Thus some remaining refinement \(\mathcal C'\) satisfies
\begin{equation}
 \#\mathcal C'\ge S_q/\mathcal Z^3
 \label{eq:inverse-refined-size}
\end{equation}
and satisfies
\[
 \frac{\#\bigl(\mathcal I_{\mathbf u}\cap
                 (\mathcal P\times\mathcal C')\bigr)}
      {n_p\#\mathcal C'}\ge\frac{\gamma_{\rm inc}}4.
\]
At least a \(\gamma_{\rm inc}/8=c_5\) fraction of its \(q\)-points consequently
have at least \(c_5 n_p\) primes \(p\) with
\((p,q)\in\mathcal I_{\mathbf u}\).
Call these points \emph{rich}, and let \(\mathcal X\) be their set of
integer points \((q,b_q)\). We have therefore fixed a sample, a refinement
\(\mathcal C'\), and slope sets \(\mathcal T_p(\mathcal C')\) with
\(\#\mathcal T_p(\mathcal C')\le K_p\). For large \(z\), their bounds are
\begin{equation}
\begin{gathered}
 |\mathcal X|\ge \frac{S_q}{\mathcal Z^4},\\
 \#\{p\in\mathcal P:(p,q)\in\mathcal I_{\mathbf u}\}\ge c_5n_p
       \quad\bigl((q,b_q)\in\mathcal X\bigr),\\
 \sum_{p\in\mathcal P}K_p\le\frac{n_pR}{L^6}.
\end{gathered}
 \label{eq:inverse-rich-points}
\end{equation}
The point-count bound follows from \eqref{eq:inverse-refined-size};
the slope sum is \eqref{eq:inverse-good-sample}. Every sampled incidence
with \(q\in\mathcal C'\) uses the slope
\(s_p(q)\in\mathcal T_p(\mathcal C')\).
The constants in these bounds are precisely those fixed in
\eqref{eq:inverse-fixed-constants}, without choosing them after the
list or after the failure set.

\medskip
\textbf{Algebraic concentration and the fixed list.}
We will find a rational line containing many rich points and show that
its intercept belongs to \(\mathcal A\). Almost all sampled incidences
on that line will then be structured, giving the contradiction.
First the small slope sum lets us construct a common polynomial equation
for the rich points.
Put
\[
 D_0=\left\lfloor\frac{R}{L^2}\right\rfloor,
 \qquad N_0=\binom{D_0+2}{2}.
\]
We construct a nonzero integer polynomial \(P(X_1,X_2)\) of total
degree at most \(D_0\) that modulo each \(p\in\mathcal P\) vanishes
identically on every line
\[
 X_2=a_p-sX_1,\qquad s\in\mathcal T_p(\mathcal C').
\]
The coefficients of the polynomial after substitution in a line give
at most \(D_0+1\) homogeneous congruences.
The lattice of coefficient vectors in \(\Z^{N_0}\) satisfying all
these congruences consequently has index at most
\[
 \prod_{p\in\mathcal P}
      p^{(D_0+1)|\mathcal T_p(\mathcal C')|}
 \le
 \exp\!\left((D_0+1)\sum_{p\in\mathcal P}K_p\log p\right).
\]
Dependencies among congruences can only decrease this index.
Applying the pigeonhole principle to a box of coefficient vectors,
and taking the difference of two vectors in the same lattice coset,
gives a nonzero \(P\) with maximum absolute coefficient \(H_P\) satisfying
\begin{equation}
 \log(2H_P)
 \ll 1+\frac1{D_0}\sum_{p\in\mathcal P}K_p\log p
 =O(R/L^4)=o(R).
 \label{eq:inverse-polynomial-height}
\end{equation}
For the final equality use \eqref{eq:inverse-good-sample},
\(n_p\ll R/\log R\), and \(\log p\asymp\log R\).
More explicitly, a coefficient box with more vectors than the
lattice index has two in the same coset; taking side length equal
to the ceiling of the index to the power \(1/N_0\) proves the stated
height bound.

All the points \((q,b_q)\) under consideration lie in \([0,S_q]^2\).
If \(P(q,b_q)\ne0\), then
\begin{equation}
 \log|P(q,b_q)|
 \le \log N_0+\log H_P+D_0\log\max(1,S_q)
 =O(R/L)=o(R).
 \label{eq:inverse-polynomial-value}
\end{equation}
At a rich point the integer \(P(q,b_q)\) is divisible by every prime
of its at least \(c_5n_p\) distinct sampled incidences.
Their product has logarithm at least
\[
 c_5n_p\log R\gg_{\xi,c_5}R.
\]
For each fixed \(\xi>0\), this contradicts
\eqref{eq:inverse-polynomial-value} unless \(P(q,b_q)=0\).
Every point of \(\mathcal X\) therefore lies on the zero set of \(P\).

\medskip
\emph{Extraction of a rational line.}
By \cref{lem:nonlinear-points}, the number of points of \(\mathcal X\)
on nonlinear irreducible components is at most
\[
 O\!\left(D_0^4(1+S_q^{2/3})\right)
 =o(S_q/\mathcal Z^4).
\]
To check this uniformly, the principal ratio is bounded by
\[
 R^4S_q^{-1/3}\mathcal Z^4
 =z^{(13\alpha-2)/3+o(1)}=o(1),
\]
using \(\alpha\le1/100\); the term without \(S_q^{2/3}\) is smaller.
There are at most \(D_0\) line components.
Discard any line containing fewer than
\[
 N_{\min}=\frac{S_q}{R\mathcal Z^6}
\]
selected rich points.
All discarded lines together contain at most
\[
 D_0N_{\min}\le\frac{S_q}{L^2\mathcal Z^6}
 =o(S_q/\mathcal Z^4)
\]
such points.
Vertical lines have at most one selected point because \(q\) uniquely
determines \(b_q\), and \(N_{\min}\to\infty\).
A line not proportional to a real line likewise has at most one real
point.  Thus a remaining nonvertical real line contains \(N\ge
N_{\min}\) selected rich points.

Two distinct lattice points on this line give it a rational slope.
It can be written
\begin{equation}
 D'b_q=A'q+B',\qquad
 D'>0,\quad \gcd(A',D')=1,\quad A',B',D'\in\Z.
 \label{eq:inverse-rich-line}
\end{equation}
For its lattice points, consecutive distinct \(q\)'s differ by an
integer multiple of \(D'\).  Their \(b_q\)'s differ by the corresponding
multiple of \(A'\).
Their coordinate ranges are at most \(S_q\), so
\begin{equation}
 D',|A'|\ll \frac{S_q}{N}\ll R\mathcal Z^6,
 \qquad
 |B'|\ll S_qR\mathcal Z^6.
 \label{eq:inverse-line-heights}
\end{equation}
The assertion for \(A'\) also holds for a horizontal line, when
\(A'=0,D'=1\).

Reduce the intercept \(B'/D'\) to \(A/D\), putting
\[
 g=\gcd(B',D'),\qquad D'=gD,\qquad B'=gA.
\]
Reduction of \eqref{eq:inverse-rich-line} modulo \(g\) gives
\(g\mid A'q\), and \(\gcd(A',g)=1\); hence
\begin{equation}
 g\mid q
 \quad\text{at every selected point of the line}.
 \label{eq:inverse-g-divides}
\end{equation}
If a prime \(\ell\mid q\) does not divide \(g\), division of
\eqref{eq:inverse-rich-line} by \(g\) modulo \(\ell\) gives
\[
 D b_q\equiv A\pmod\ell.
\]
Because \(b_q\equiv a_\ell\pmod\ell\), the prime \(\ell\) aligns with
the intercept.  This congruence also implies \(\ell\nmid D\), since
\(\gcd(A,D)=1\).
All nonaligned prime factors of \(q\) therefore divide \(g\).
As \(q\) is squarefree and \(g\mid q\), their product divides \(g\),
and
\begin{equation}
 Dq_E(A/D)\le Dg=D'\ll R\mathcal Z^6.
 \label{eq:inverse-line-structure}
\end{equation}
Neither \(D\) nor \(D'\) has been assumed squarefree.
In particular, primes common to \(D\) and \(g\) cause no difficulty in
\eqref{eq:inverse-line-structure}.
The height bounds in \eqref{eq:inverse-list} follow from
\eqref{eq:inverse-line-heights} with ample room in the additional
\(\mathcal Z^4\) factors.
It remains to establish the alignment-count condition for this
intercept.

\medskip
\emph{The intercept belongs to the list.}
We use the pair-counting principle of Gallagher's larger sieve
\cite{Gallagher1971}; see \cite[Section~2]{HelfgottVenkatesh2009} for
the difference-product argument.
Ignore the isolated-bin primes dividing \(D'\).
There are at most
\[
 \frac{\log D'}{\log R}=O(1)
\]
of them by \eqref{eq:inverse-line-heights}.
Let \(\mathcal P_{\rm un}\) consist of the other primes of
\(\mathcal P\) that do not align with \(A/D=B'/D'\).
For \(p\in\mathcal P_{\rm un}\), substitution of the line equation
into the incidence equation gives
\[
 (A'/D'+s)q\equiv a_p-B'/D'\pmod p,
 \qquad s\in\mathcal T_p(\mathcal C').
 \]
The right side is nonzero because \(p\) is unaligned.
Each allowed \(s\) therefore admits either no solution or exactly one
class of \(q\bmod p\).
Thus those of the \(N\) selected points on the line that satisfy the
sampled constraints for \(p\) occupy at most \(K_p\) classes of \(q\bmod p\).
Let their number be \(N_p\), and let \(n_{p,c}\) denote the numbers
in these occupied classes. Here \(N_p\) counts all tested points on the
line, so it dominates the number of pairs
\((p,q)\in\mathcal I_{\mathbf u}\) there; the latter pairs also require
the parent to be nonstructured.

For two distinct selected points, \(q_1-q_2\ne0\) and
\(|q_1-q_2|\le S_q\).
The number of isolated-bin primes dividing this difference is at
most
\[
 \frac{\log S_q}{\log R}=O_{\alpha_{\min}}(1).
\]
Counting ordered equal-class pairs, with the diagonal separated,
therefore gives
\[
 \sum_{p\in\mathcal P_{\rm un}}\sum_c n_{p,c}^2
 \le C_{\alpha_{\min}}N^2+
       \sum_{p\in\mathcal P_{\rm un}}N_p.
\]
Cauchy's inequality, first over all occupied classes and primes,
yields
\begin{equation}
 \left(\sum_{p\in\mathcal P_{\rm un}}N_p\right)^2
 \le
 \left(\sum_{p\in\mathcal P}K_p\right)
 \left(C_{\alpha_{\min}}N^2+
             \sum_{p\in\mathcal P_{\rm un}}N_p\right).
 \label{eq:inverse-collisions}
\end{equation}
Use \(\sum_p N_p\le n_pN\), the good-sample bound, and
\(N\ge S_q/(R\mathcal Z^6)\).
After division by \(n_p^2N^2\), the right side of
\eqref{eq:inverse-collisions} is
\begin{equation}
 O\!\left(\frac{R}{n_pL^6}+\frac{R}{NL^6}\right)
 =O_\xi(L^{-5})+o(1)=o(1).
 \label{eq:inverse-collision-normalization}
\end{equation}
For the first term use \eqref{eq:inverse-bin-counts} and
\(\log R\le L\); for the second use
\(R^2\mathcal Z^6/S_q=z^{3\alpha-2+o(1)}=o(1)\).
Thus unaligned primes supply only \(o(n_pN)\) tested pairs on the line,
and hence at most this many sampled incidences.

Each of its \(N\) selected points has at least \(c_5n_p\)
sampled incidences by \eqref{eq:inverse-rich-points}.
The ignored divisors of \(D'\) supply at most \(O(N)\).
The remaining incidences force at least
\((c_5-o(1))n_p\) primes of \(\mathcal P\) to align with the intercept.
Since \(c_p=c_5/4\), it satisfies the last condition in
\eqref{eq:inverse-list}, and hence \(A/D\in\mathcal A\).
At every selected point of this line,
\eqref{eq:inverse-line-structure} implies the required bound
\(Dq_E(A/D)\le R\mathcal Z^{10}\) for large \(z\).
Every tested pair on the line whose isolated prime aligns is therefore
structured. Since all members of \(\mathcal I_{\mathbf u}\) are
nonstructured, the sampled incidences must consequently come from the
unaligned primes or the \(O(1)\) ignored primes; those
contribute only \(o(n_pN)+O(N)\), contradicting their lower bound
\(c_5n_pN\).

This contradiction proves the exceptional-fraction assertion.
The constants in \eqref{eq:inverse-fixed-constants} and the list bound
were fixed before \(\xi\).
Only interval freezing required an upper restriction
\(\xi\le\xi_{\rm inv}\); all losses that depend on a fixed positive
\(\xi\), including prime counts, polynomial divisibility, and
\eqref{eq:inverse-collision-normalization}, are absorbed by taking
\(z\) sufficiently large afterwards.
This completes the prescribed order of choices.
\end{proof}
